\input{../tex-inputs/latex-leseansicht-vorspann}

\section[Gustav und Max Schwarzkopf an Arthur Schnitzler, 6. 5. 1904]{L04315 Gustav und Max Schwarzkopf an Arthur Schnitzler, 6. 5. 1904}
\nopagebreak\mylabel{L04315v}
\rehead{ }\normalsize\beginnumbering\briefempfaengerindex{Schnitzler, Arthur@\textsc{Schnitzler, Arthur}!zzzSchwarzkopf, Max@\emph{von Max Schwarzkopf}!1904-05-061@{6. 5. 1904}|(be}\briefempfaengerindex{Schnitzler, Arthur@\textsc{Schnitzler, Arthur}!zzzSchwarzkopf, Gustav@\emph{von Gustav Schwarzkopf}!1904-05-061@{6. 5. 1904}|(be}
\toendnotes[C]{\smallbreak\pagebreak[2]}
\correspDesc{Versand  durch Gustav Schwarzkopf, Max Schwarzkopf am 6. 5. 1904 in Wien
\newline{}Übermittlung  am 8. 5. 1904 in Wien
\newline{}Zustellung  am 9. 5. 1904 in Neapel
\newline{}Erhalt  durch Arthur Schnitzler am 9. 5. 1904 in Neapel}\toendnotes[C]{\smallbreak}
\Standort{CUL, Schnitzler, B 96a.}
\physDesc{Postkarte, 401 Zeichen
\newline{}Handschrift Gustav Schwarzkopf: schwarze Tinte, deutsche Kurrent
\newline{}Handschrift Max Schwarzkopf: schwarze Tinte
\newline{}Versand: Stempel: »\nobreak{}\oindex{I., Innere Stadt@\textbf{I., Innere Stadt}, \emph{Verwaltungsgebiet}|pwk}Wien 1/1, 8. 5. \textcolor{gray}{04}, \textcolor{gray}{6}\nobreak{}«.  }\toendnotes[C]{\smallbreak}\pstart{}{\pb}\textsc{Sig. Arthur Schnitzler}\pend{}\pstart{}\textsc{Neapel\oindex{Neapel@\textbf{Neapel}|pw}}\pend{}\pstart{}\textsc{poste restante}\pend{}{\bigskip}\vspace{1em}
\pstart
           \raggedleft{}{\pb}I. Tiefer Graben 23. Th. 6\oindex{Wien@\textbf{Wien}!I., Innere Stadt@\textbf{I., Innere Stadt}!Tiefer Graben 23@\textbf{Tiefer Graben 23}, \emph{Wohngebäude}|pw}.\pend
           
\pstart
           \raggedleft{}\textsc{Wien\oindex{Wien@\textbf{Wien}, \emph{Verwaltungsgebiet}|pw}}6. V 04\pend
           \vspace{0.5em}
\pstart
           Lieber Arthur! Iſt das bei Ihnen befindliche Exemplar »\uline{der reine Tor\pwindex{Schwarzkopf, Max 12.\,6.\,1857 Wien – 14.\,4.\,1928 ebd.@\textsc{Schwarzkopf, Max} (12.\,6.\,1857 Wien – 14.\,4.\,1928 ebd.), \emph{Rechtsanwalt}!reine Tor. Gesellschaftsstück in vier Akten@\strich\emph{Der reine Tor. Gesellschaftsstück in vier Akten}|pw}}« für mich erreichbar, d.h. liegt es vielleicht in einem Ihrer Bücherschränke,
               daß ich es abholen kö{\geminationn}te? (Prag\oindex{Prag@\textbf{Prag}, \emph{Land}|pw}\orgindex{Neues Deutsches Theater@Neues Deutsches Theater|pwv} verlangt nämlich 2. Exemplare und Max beſitzt nur
               eines.) Bitte mir um eine kurze Verſtändigung.\pend
           
\pstart
           Ihre Frau\pwindex{Schnitzler, Olga 17.\,1.\,1882 Wien – 13.\,1.\,1970 Lugano@\textsc{Schnitzler, Olga} (17.\,1.\,1882 Wien – 13.\,1.\,1970 Lugano), \emph{Schauspielerin, Sängerin}|pwv} und Sie
               herzlichſt grüßend{\\[\baselineskip]}Ihr{\\[\baselineskip]}\spacefill\mbox{Gustav.}{\\[\baselineskip]}\spacefill\mbox{{[}hs. Schwarzkopf:{]} MaxS}\pend
           \leftskip=0em{}\selectlanguage{ngerman}\endnumbering\briefempfaengerindex{Schnitzler, Arthur@\textsc{Schnitzler, Arthur}!zzzSchwarzkopf, Max@\emph{von Max Schwarzkopf}!1904-05-061@{6. 5. 1904}|)be}\briefempfaengerindex{Schnitzler, Arthur@\textsc{Schnitzler, Arthur}!zzzSchwarzkopf, Gustav@\emph{von Gustav Schwarzkopf}!1904-05-061@{6. 5. 1904}|)be}\mylabel{L04315h}
\begin{anhang}
\end{anhang}\newcommand{\dateiname}{L04315}\newcommand{\titel}{Gustav und Max Schwarzkopf an Arthur Schnitzler, 6. 5. 1904}\newcommand{\editorInnen}{Herausgegeben von Jahnke, SelmaMüller, Martin Anton}\input{../tex-inputs/latex-leseansicht-abspann}
